\documentclass[../../../../main.tex]{subfiles}
\begin{document}

\begin{flushright}
\footnotesize\textit{【自動文字起こし・要確認】}
\end{flushright}

[解]\\
(1) $b, c \in \mathbb{R}$ であるから、この時、$\beta = \overline{\alpha}$ となるので、
\[
\begin{cases}
\alpha + \beta = -2b \quad (\neq 0) \\
\alpha - \beta = 2id \quad (d \in \mathbb{R})
\end{cases}
\]
とおける。したがって、$\gamma = v + iw$ とすると ($v, w \in \mathbb{R}$)\\
$\gamma = \frac{v}{2b}(\alpha + \beta) + \frac{w}{2d}(\alpha - \beta) = \left( -\frac{v}{2b} + \frac{w}{2d} \right) \alpha + \left( -\frac{v}{2b} - \frac{w}{2d} \right) \beta$\\
と表せる。$\blacksquare$

(2) (1)から、$b^2 - c \ge 0$ 又は $b = 0$ が必要である。\\
ア) $b^2 - c \ge 0$ の時\\
$\alpha, \beta$ は実数である。$\alpha = \beta = 0 \iff b = c = 0$ に注意する。
\begin{itemize}
\item $b = c = 0$ の時、$f(x) = 5 \neq 0$ より $(\ast)$ をみたす。
\item otherwise、$t\alpha + u\beta$ は任意の実数値をとるから、
常に $f(x) \neq 0 \iff (b-1)^2 - 5 + c < 0 \iff c < -(b-1)^2 + 5$
\end{itemize}
イ) $b = 0$ かつ $b^2 - c < 0$ の時\\
$\alpha, \beta$ は純虚数であるから、$t\alpha + u\beta$ は任意の $ki$ ($k \in \mathbb{R}$) の形の値をとる。\\
$f(ki) = (-k^2 - c + 5) - 2ki = 0 \iff k = 0 \text{ かつ } c = 5$\\
より、$c \neq 5$ であれば $(\ast)$ をみたす。

ア)、イ) より、もとめるのは、\\
"$b^2 - c \ge 0$ かつ $c < -(b-1)^2 + 5$"\\
or\\
"$b = 0$ かつ $b^2 - c < 0$ かつ $c \neq 5$"\\
で、図示すると右図斜線部(境界は実線のみ含む)

\begin{tikzpicture}
\draw[->] (-3,0) -- (4,0) node[right] {$b$};
\draw[->] (0,-2) -- (0,7) node[above] {$c$};
\fill[pattern=north east lines] (0,0) parabola (2,4) -- (2,4) parabola bend (1,5) (-1,1) -- (-1,1) parabola bend (0,0) (0,0);
\draw[thick] (0,0) parabola (2,4);
\draw[thick] (0,0) parabola (-1,1);
\draw[dashed] (-1,1) parabola bend (1,5) (3,-3);
\draw[thick] (0, 0.05) -- (0, 4.95);
\draw[thick] (0, 5.05) -- (0, 7);
\draw[fill=white] (0,5) circle (0.05);
\node at (0.2, 5) {$5$};
\node at (1.5, 4.5) {$c = b^2$};
\node at (2.5, 6) {$c = -(b-1)^2 + 5$};
\end{tikzpicture}

\end{document}
